\providecommand{\C}{\mathbb{C}}
\providecommand{\Z}{\mathbb{Z}}
\providecommand{\N}{\mathbb{N}}
\providecommand{\R}{\mathbb{R}}
\providecommand{\Q}{\mathbb{Q}}
\providecommand{\gl}{\mathfrak{gl}}
\providecommand{\bowtie}{\bowtie}
\providecommand{\CoHA}{\mathrm{CoHA}}
\providecommand{\Hilb}{\mathrm{Hilb}}
\providecommand{\Mbps}{\mathcal{M}^{+}_{\mathrm{eff}}}
\providecommand{\Meff}{\mathcal{M}_{\mathrm{eff}}}
\providecommand{\Yplus}{Y^{+}}
\providecommand{\Yminus}{Y^{-}}
\providecommand{\Yzero}{Y^{0}}
\providecommand{\Heq}{H^{*}_{\mathrm{eq}}}
\providecommand{\kchir}{\kappa_{\mathrm{ch}}}
\providecommand{\kcat}{\kappa_{\mathrm{cat}}}
\providecommand{\kBKM}{\kappa_{\mathrm{BKM}}}
\providecommand{\kfib}{\kappa_{\mathrm{fiber}}}
\providecommand{\bps}{\mathrm{BPS}}
\providecommand{\red}{\mathrm{red}}
\providecommand{\ClaimStatusTheorem}{\textsc{[theorem, cited]}}
\providecommand{\ClaimStatusDefinition}{\textsc{[definition]}}
\providecommand{\ClaimStatusDerivation}{\textsc{[first-principles derivable]}}
\providecommand{\ClaimStatusConjectured}{\textsc{[conjectural]}}
\providecommand{\ClaimStatusOpen}{\textsc{[open]}}

\section*{Wave-17 A1/20: the universal grammar
         $\Yplus(X) = \Heq(\Mbps(X))$}
\label{wn:sec:universal-Yplus-grammar}

The positive half is not a family of tricks; it is a single grammar
read across four equivariance strata.  The grammar instructs: given a
Calabi--Yau datum $X$ of dimension $d$, select the positive effective
moduli stack $\Mbps(X)$ carrying the Kontsevich--Soibelman critical
structure, and form its equivariant cohomology under the symmetry
that survives compactness.  What follows formalises this, then walks
it across toric CY$_{3}$, non-toric compact CY$_{3}$, orbifold, and
lattice-polarised strata, closing on $K3 \times E$, Drinfeld doubling,
and the $E_{d}$-hFA selection rule.

\subsection*{Attack-heal 1 --- the universal grammar}

\ClaimStatusDefinition\ \emph{Positive effective moduli.}  For a
Calabi--Yau $\infty$-category $X$ of dimension $d$, the \emph{positive
effective moduli stack} $\Mbps(X) \subset \Meff(X)$ is the substack of
objects of positive charge $\gamma \in \Gamma^{+}(X)$ with respect to
the Kontsevich--Soibelman central charge, equipped with the critical
structure $(W, \phi_{W})$ coming from the minimal Maurer--Cartan
presentation of the Calabi--Yau shift (Kontsevich--Soibelman 2008,
\texttt{arXiv:0811.2435}; Van den Bergh 2004, \emph{NCCR} construction).

\ClaimStatusDefinition\ \emph{Universal grammar.}
\[
   \Yplus(X) \;=\; \Heq\bigl(\Mbps(X), \phi_{W}\bigr)
   \;=\; \Heq\bigl(\Meff(X), \phi_{W}\bigr)\bigr|_{\gamma \in \Gamma^{+}(X)}.
\]
The two writings agree: in the first the positive cone is absorbed
into the moduli; in the second (matching Wave-17 A2/20) the
positive cone is cut by the equivariance.  The equivariance $\Heq$
is fixed by the next stratum.  The
multiplication is the Kontsevich--Soibelman correspondence product on
the Hall-type stack $\mathcal{M}_{\gamma_{1}} \times \mathcal{M}_{\gamma_{2}}
\xleftarrow{} \mathcal{M}_{\gamma_{1} \to \gamma_{1} + \gamma_{2}}
\xrightarrow{} \mathcal{M}_{\gamma_{1} + \gamma_{2}}$.  The grammar
reads as a single formula across $X$; what changes is which
equivariance survives compactness.

\subsection*{Attack-heal 2 --- the four equivariance strata}

\ClaimStatusDerivation\ \emph{Stratification.}  Four strata exhaust the
programme's CY inputs:

\begin{enumerate}
\item \emph{Toric CY$_{3}$.}  $X$ admits a Calabi--Yau torus
$T = (\C^{\times})^{3}$ with $t_{1} t_{2} t_{3} = 1$.  $\Heq =
H^{*}_{T}$, and localisation reduces to $T$-fixed quiver
representations (Schiffmann--Vasserot 2013, \texttt{arXiv:1202.2756},
Thm 1.1).  This is the $\Yplus(\C^{3}) = \CoHA(\C^{3})$ benchmark:
fixed loci are Young diagrams, characters are MacMahon.

\item \emph{Non-toric compact CY$_{3}$.}  No non-trivial torus acts;
one replaces $T$-equivariance by the \emph{reduced virtual class plus
residual $\C^{\times}$}.  The virtual cotangent complex is perfect of
virtual dimension $0$ (Pandharipande--Thomas 2014,
\texttt{arXiv:0904.2891}, on stable pairs;
Maulik--Nekrasov--Okounkov--Pandharipande 2006
\texttt{arXiv:math/0312059} for the reduced theory on surfaces
fibered over $\C^{\times}$).  The residual $\C^{\times}$ is the
holomorphic-symplectic scaling on the elliptic / abelian direction or,
when the CY$_{3}$ is a $K3$-fibration, the $\C^{\times}$ acting on the
base.

\item \emph{Orbifold / finite-symplectic.}  When $X = [Y/G]$ with $G$
finite symplectic, $\Heq = H^{*}_{G}$ acts through the Nikulin--Mukai
$M_{23}$-subgroup classification.  The positive effective stack
fibres over the discrete fixed strata $Y^{G}$, and $\Yplus([Y/G])$
receives a Davison--Meinhardt 2020 (\emph{Invent. Math.}\ 221) BPS Lie
algebra by descent through the orbifold cotangent complex.

\item \emph{Lattice-polarised.}  When $X$ is polarised by a lattice
$L \subset H^{2}(X, \Z)$, the automorphic moduli stack (period domain
$\mathcal{A}_{2} = \mathrm{Sp}(4, \Z) \backslash \mathbb{H}_{2}$ for
$K3$ type) carries a \emph{period-domain equivariance}: $\Heq =
H^{*}_{\mathrm{O}(L)}$, equivariant under the orthogonal group of
the lattice.  The character of $\Yplus(X)$ is then an automorphic form
on $\mathcal{A}_{L}$, matching the Gritsenko--Nikulin 1998
(\emph{Amer. J. Math.}\ 119) Borcherds lifts that define $\kBKM$.
\end{enumerate}

In all four strata, the multiplicative structure is the KS
correspondence product; only the coefficient ring
$\mathbb{F} = \Heq(\mathrm{pt})$ shifts.

\subsection*{Attack-heal 3 --- $K3 \times E$ via Pandharipande--Oberdieck reduced DT}

\ClaimStatusTheorem\ (Oberdieck--Pixton 2017 \texttt{arXiv:1706.10100},
\emph{Invent. Math.}\ 222 (2020), primitive $K3$ class.)
\[
  Z^{\red}_{DT}\bigl(K3 \times E; \gamma \text{ primitive on } K3\bigr)(q, t, p)
   \;=\; -\,\frac{C}{\Phi_{10}(q, t, p)}.
\]

\ClaimStatusDerivation\ \emph{Positive effective geometry for
$K3 \times E$.}  Stratum 2 applies: $K3$ is non-toric; $E$ supplies the
surviving $\C^{\times}_{E}$ via its translation automorphism group.
The positive effective moduli is the Pandharipande--Oberdieck reduced
DT stack,
\[
  \Mbps(K3 \times E) \;=\;
    \mathcal{M}^{\red, \bps}\!\bigl(K3 \times E;\, \gamma\bigr)_{\gamma \text{ $K3$-primitive}},
\]
with reduced virtual dimension $0$ (the $K3$-holomorphic-symplectic
direction is excised, per MNOP reduced theory), carrying the residual
$\C^{\times}_{E}$-action from $E$-translation.  The universal grammar
then gives
\[
  \Yplus(K3 \times E) \;=\; H^{*}_{\C^{\times}_{E}}\bigl(\mathcal{M}^{\red, \bps}(K3 \times E), \phi_{W}\bigr).
\]
Restricting to $K3$-primitive charges, the $\C^{\times}_{E}$-equivariant
character reproduces $\Phi_{10}^{-1}$ via Oberdieck--Pixton 2017, which
locks $\kBKM(\Phi_{10}) = c_{1}(0)/2$ to the Gritsenko $\Delta_{5}$
weight (Lorgat 2020, April; $\kBKM = 5$).  The imprimitive
multiple-cover contribution sits outside the stated identity and is
the $\Yplus$-side opening for Maulik--Toda 2018 refinement.

\ClaimStatusOpen\ \emph{Lie-algebra isomorphism
$\mathfrak{g}_{\bps}(K3 \times E) \cong \mathfrak{g}_{\Delta_{5}}$.}
Charge-by-charge dimensions match via KS wall-crossing (Attack-heal 3
of \texttt{CoHA\_to\_W\_infty\_treatise.tex}); the bracket-level
identification remains open (non-adjacent contractions, per AP-CY34).

\subsection*{Attack-heal 4 --- Drinfeld double commutes with the stratum selector}

\ClaimStatusDefinition\ \emph{Drinfeld double.}
Given $\Yplus(X)$ with a non-degenerate Hopf pairing
$\langle -, - \rangle: \Yplus(X) \otimes \Yminus(X) \to \Yzero(X)$,
the Drinfeld double is
\[
  G(X) \;=\; \Yplus(X) \bowtie \Yzero(X) \bowtie \Yminus(X).
\]

\ClaimStatusDerivation\ \emph{Pairing source.}
In the compact stratum (2, 4), the pairing is Serre duality on
$D^{b}(X)$: $\mathrm{Ext}^{d}_{X}(E, F)^{\vee} \cong \mathrm{Ext}^{0}_{X}(F, E)$
for $E, F$ in the perfect Hall-stack, descending to a non-degenerate
pairing on $\Heq$-cohomology via the Davison--Meinhardt 2020 PBW
integration map.  In the toric stratum (1), the pairing is the
equivariant residue on torus-fixed loci; the two agree on the
intersection by the Atiyah--Bott fixed-point formula.

\ClaimStatusDerivation\ \emph{Universality.}
The doubling constructor commutes with the stratum selector.  At
$\C^{3}$, $G(\C^{3}) = Y(\widehat{\gl}_{1})$ (Tsymbaliuk 2017
\texttt{arXiv:1703.04551}).  At the conifold, $G(\mathrm{cfd}) = Y^{\mathrm{cfd}}$
(Szendroi 2008 \texttt{arXiv:0708.4384} NCCR quiver).  At $K3 \times E$,
$G(K3 \times E) \supset \mathfrak{g}_{\Delta_{5}}$ conjecturally,
with Cartan series $\psi(z)$ automorphic on $\mathcal{A}_{2}$.  Six
routes to $G(K3 \times E)$ are six different \emph{constructions},
not six $\Phi$-applications.

\subsection*{Attack-heal 5 --- two-stage coherence: $E_{d}$-hFA selection}

\ClaimStatusDerivation\ \emph{Stage-2 selection.}
In the Vol III two-stage $E_{d}$-hFA framework, \emph{Stage 1} is the
raw chiral factorisation $\mathrm{Fact}_{E_{d}}$ on the Ran space;
\emph{Stage 2} is $\mathrm{Sp}_{\Sigma_{d - 1}, C}$, the
$\Sigma_{d - 1}$-stratified specialisation along a CY curve or surface
slice $C$.  Stage 2 is what \emph{selects} the positive effective
geometry: the cone $\Gamma^{+}(X)$ is the image of the KS central
charge read through Stage 2, and the equivariance stratum is the
monodromy of the specialisation.

\ClaimStatusDerivation\ \emph{Coherence.}  Writing
$\Phi: \mathrm{CY}\text{-cat}_{d} \to \mathrm{ChirAlg}$ for the
CY-to-chiral functor, the universal grammar gives a square
\[
  \Phi(X) \xhookrightarrow{\text{Stage 1}} \mathrm{Fact}_{E_{d}}(X)
  \xrightarrow{\mathrm{Sp}_{\Sigma_{d-1}, C}} \Yplus(X) =
  \Heq(\Mbps(X), \phi_{W}),
\]
the horizontal composite agreeing with the Hodge-supertrace
identification $\kchir(A_{X}) = \sum_{q} (-1)^{q} h^{0, q}(X)$ on
compact CY$_{d}$.

At $d = 3$, Stage 2 picks out the reduced virtual cycle and the
residual $\C^{\times}_{E}$; at $d = 2$, it picks out the $K3$ Mukai
lattice and its $\mathrm{O}(\Lambda_{K3})$-equivariance.  The grammar
is one formula; the coherence is that Stage 2 always reads it in the
correct stratum.

\subsection*{Status}

The grammar is written.  Strata 1 (Schiffmann--Vasserot) and 3
(Nikulin--Mukai orbifold) are theorems.  Stratum 2 at $K3 \times E$ is
a definition plus the Oberdieck--Pixton primitive-class identification,
with the bracket-level $\mathfrak{g}_{\bps}(K3 \times E) \cong
\mathfrak{g}_{\Delta_{5}}$ isomorphism open.  Stratum 4 is programmatic,
anchored to Gritsenko--Nikulin 1998.  The Drinfeld doubling commutes
with the stratum selector; the two-stage $E_{d}$-hFA provides the
selection rule.  Open: bracket-level compact-CY$_{3}$ identification,
the automorphic $\psi(z)$ on $\mathcal{A}_{2}$, and the lattice-polarised
stratum character formula for a general polarisation $L \ne U \oplus U
\oplus U \oplus E_{8}(-1)^{\oplus 2}$.
